% Generated draft: requires mathematical and source review.
\documentclass[11pt]{article}
\usepackage[margin=25mm]{geometry}
\usepackage{fontspec}
\setmainfont{Latin Modern Roman}
\usepackage{amsmath,amssymb,mathtools,unicode-math}
\setmathfont{Latin Modern Math}
\usepackage{enumitem,needspace,xurl,hyperref}
\setlength{\parindent}{0pt}
\setlength{\parskip}{5pt}
\setlength{\emergencystretch}{3em}
\newcommand{\N}{\mathbb{N}}
\newcommand{\Z}{\mathbb{Z}}
\newcommand{\Q}{\mathbb{Q}}
\newcommand{\R}{\mathbb{R}}
\newcommand{\C}{\mathbb{C}}
\newcommand{\F}{\mathbb{F}}
\newcommand{\A}{\mathbb{A}}
\newcommand{\PP}{\mathbb{P}}
\newcommand{\E}{\mathbb{E}}
\newcommand{\eps}{\varepsilon}
\newcommand{\dd}{\,\mathrm{d}}
\begin{document}
\section*{UnsolvedMath ID: 700505}
\textbf{Source problem code:} catalogue entry
\subsection{Definitions and mathematical statement}
A quasivariety of groups is a class of groups axiomatized by universally quantified quasi-identities of the form
\[
\bigl(w_1(x_1,\ldots,x_n)=1\ \land\cdots\land\ w_m(x_1,\ldots,x_n)=1\bigr)\ \Longrightarrow\ w(x_1,\ldots,x_n)=1,
\]
where the expressions are group words and \(m\geq 0\). A nontrivial identity means an identity \(v(x_1,\ldots,x_n)=1\), holding for every substitution in every group in the class, where \(n\geq 1\) and \(v\) is a nonidentity element of the free group on \(x_1,\ldots,x_n\).

Let
\[
\mathcal S=\{\mathcal Q:\mathcal Q\text{ is a quasivariety of groups satisfying some nontrivial identity}\}.
\]
The identity may depend on \(\mathcal Q\). Under the usual extension-product convention, multiplication is
\[
\mathcal Q\mathcal R
=\{G:\text{there exists }N\trianglelefteq G\text{ with }N\in\mathcal Q\text{ and }G/N\in\mathcal R\}.
\]
The supplied source asserts that these quasivarieties form a semigroup under multiplication and asks whether this semigroup is free.

Explicitly, freeness as a semigroup would mean that there exists a set \(B\subseteq\mathcal S\) such that every element of \(\mathcal S\) has a unique expression as a nonempty finite product of elements of \(B\). Thus, for \(r,s\geq 1\) and \(\mathcal Q_i,\mathcal R_j\in B\),
\[
\mathcal Q_1\cdots\mathcal Q_r=\mathcal R_1\cdots\mathcal R_s
\quad\Longrightarrow\quad
r=s\ \text{ and }\ \mathcal Q_i=\mathcal R_i\ (1\leq i\leq r).
\]
\textbf{Unresolved conventions.} The source does not define multiplication; the displayed product makes the customary extension interpretation explicit. Also, as written, \(\mathcal S\) includes the quasivariety \(\mathcal T\) consisting only of trivial groups, which is an identity element for this product. The source does not specify whether its freeness question excludes \(\mathcal T\), or instead uses freeness with an identity adjoined, allowing the empty product to represent \(\mathcal T\). These conventions must be clarified before fixing the intended freeness assertion; the nonempty-product formulation above states the literal free-semigroup reading.
\subsection{Short English statement}
A. Yu. Olshanskii asks whether quasivarieties of groups that each satisfy some nontrivial identity form a free semigroup under multiplication. This is a unique factorization question for products of quasivarieties. The supplied wording leaves the multiplication convention and the treatment of the trivial quasivariety as an identity element unresolved.
\subsection{Sources}
\begin{itemize}
\item Local UnsolvedMath registry: supplied complete source record, permanent ID \(700505\), exact source problem code \texttt{catalogue entry}, titled \emph{Kourovka Notebook Problem 8.55}. The record labels the problem \texttt{open}; this is a supplied registry status, not an independently verified mathematical assessment. The record was created and updated on \texttt{2026-10-16T14:00:00Z} and carries upstream revision \texttt{d371ee917ac9647fbb06cd6fc529deec6f3dda82}.
\item \emph{The Kourovka Notebook (Unsolved Problems in Group Theory)}, original Problem 8.55, proposed by A. Yu. Olshanskii in the 8th issue (1982). The supplied background identifies the 21st issue edition (2026), \texttt{arXiv:1401.0300v48}.
\item The supplied collection description attributes editing of the \emph{Kourovka Notebook}, Issues 1--21 (1965--2026), to E. I. Khukhro and V. D. Mazurov. The cited edition and attribution are supplied source information; the external text has not been independently consulted.
\end{itemize}
\end{document}
